\documentclass[12pt]{article}
\usepackage[utf8]{inputenc}
\usepackage[T1]{fontenc}
\usepackage{hyperref}
\title{Choapan Zapotec Language Resources for K-12 Curriculum Development}
\author{Compiled by Claude Sonnet 5.5 (claude-sonnet-5-5) for Andr\'as Kornai}
\date{}
\begin{document}
\maketitle

\section{Introduction}
Choapan Zapotec (ISO 639-3 zpc; Glottocode choa1237) is a Northern Zapotec language of the Choap\'am region of Oaxaca, Mexico, with about 12,000 speakers. It has three tone levels. Documentation exists, mostly from linguists' fieldwork. I found no published school materials.

\section{Primary Educational Resources}
None found. A 2012 fieldwork report by Erin Donnelly describes work on a dictionary she planned to publish through INALI. I could not confirm that it was published.
\begin{itemize}
\item Documentation of Choapan Zapotec, UC Berkeley CLACS: \url{https://clacs.berkeley.edu/node/1097}
\end{itemize}

\section{Linguistic and Reference Materials}
\begin{itemize}
\item Wikipedia, ``Cho\'apam Zapotec'' (phonology and tone): \url{https://en.wikipedia.org/wiki/Cho%C3%A1pam_Zapotec}
\item Omniglot, Choapan Zapotec (alphabet table, stress note, short text with Spanish translation): \url{https://omniglot.com/writing/zapotec-choapam.htm}
\end{itemize}

\section{Religious and Community Resources}
\begin{itemize}
\item Global Recordings Network, Zapotec, Choapan (Bible audio and evangelistic messages; it also links to Scripture Earth and Faith Comes By Hearing material): \url{https://globalrecordings.net/en/language/zpc}
\end{itemize}

\section{Digital and Audio Resources}
\begin{itemize}
\item AILLA, Northern Zapotec Collection of Erin Donnelly (interviews, oral histories, grammatical elicitation; the archive says the collection is being reorganized and asks users not to cite it): \url{https://ailla-legacy.lib.utexas.edu/islandora/object/ailla%3A124501}
\item AILLA, Choapan Northern Zapoteco (elicitation and narrative recordings from 1995 and 1998, from Terrence Kaufman's documentation project): \url{https://islandora-ailla.lib.utexas.edu/islandora/object/ailla%3A135302}
\item AILLA, Narratives in Choapan Northern Zapoteco: \url{https://islandora-ailla.lib.utexas.edu/islandora/object/ailla%3A135307}
\end{itemize}

\section{Curriculum Development Context}
Not all young people speak the language as a first language, so a curriculum is relevant to its maintenance. The language is not known to be taught in schools. Texts and recordings in the AILLA archives could serve as source material, subject to the archive's access terms. The alphabet on Omniglot is a possible orthography reference, but it should be checked with community members and INALI. Primers, textbooks and a dictionary would have to be developed or located.

\end{document}